\section{字符级语言模型}

\subsection{自回归生成}

RNN 可以训练为\textbf{字符级语言模型}：给定前面的字符，预测下一个字符。这本质上是一个逐时间步的分类问题。

\begin{figure}[H]
\centering
\includegraphics[width=0.95\textwidth]{slides-images/slide-022.jpg}
\caption{字符级语言模型：输入字符的 one-hot 编码，隐藏层更新状态，输出层产生下一个字符的概率分布\protect\footnotemark}
\end{figure}
\footnotetext{Slides 第23页。}

\begin{importantbox}{语言模型的训练与推理}
\textbf{训练时}：
\begin{itemize}
    \item 输入真实序列，目标是下一个字符
    \item 每个时间步计算交叉熵损失
    \item 无需人工标注——文本本身就是监督信号（\textbf{自监督学习}）
\end{itemize}

\textbf{推理时（生成）}：
\begin{itemize}
    \item 给定一个起始字符（或起始 token）
    \item 从输出概率分布中\textbf{采样}一个字符
    \item 将采样结果作为下一步的输入
    \item 重复直到生成结束 token 或达到最大长度
\end{itemize}
\end{importantbox}

\subsection{Embedding 层}

在实际实现中，通常不直接使用 one-hot 编码作为输入，而是使用\textbf{Embedding 层}：一个可学习的查找表，将每个离散 token 映射为一个稠密向量。

\begin{knowledgebox}{One-hot vs Embedding}
\begin{itemize}
    \item \textbf{One-hot}：稀疏、高维（维度 = 词表大小），不包含语义信息
    \item \textbf{Embedding}：稠密、低维（通常 64--512 维），通过训练自动学习语义关系
    \item 数学上，用 one-hot 向量做矩阵乘法等价于从 Embedding 矩阵中取出对应行
    \item Embedding 在优化上更高效，且能学到有意义的表示
\end{itemize}
\end{knowledgebox}

\subsection{采样策略}

\begin{warningbox}{贪心解码的问题}
如果每次都选择概率最大的字符（贪心解码），模型会反复生成相同的序列。实际做法是：
\begin{itemize}
    \item \textbf{随机采样}：按概率分布采样
    \item \textbf{温度调节}：$p_i = \frac{e^{z_i / T}}{\sum_j e^{z_j / T}}$，温度 $T$ 控制分布的``尖锐度''
    \item \textbf{Beam Search}：同时维护多个候选序列，选择整体概率最高的
    \item \textbf{Top-k / Top-p 采样}：只从概率最高的 $k$ 个（或累积概率达到 $p$ 的）token 中采样
\end{itemize}
\end{warningbox}

\subsection{生成效果展示}

用仅 112 行 Python 代码实现的字符级 RNN，训练在莎士比亚十四行诗上，可以生成风格相似的文本。随着训练的进行，生成质量逐步提升：

\begin{figure}[H]
\centering
\includegraphics[width=0.95\textwidth]{slides-images/slide-025.jpg}
\caption{训练不同阶段的生成结果：从随机乱码逐步变为有结构的、类似莎士比亚风格的文本\protect\footnotemark}
\end{figure}
\footnotetext{Slides 第26页。}

\begin{figure}[H]
\centering
\includegraphics[width=0.95\textwidth]{slides-images/slide-026.jpg}
\caption{在 Linux 源代码上训练的字符级 RNN 生成的 C 代码：看起来结构合理，虽然无法编译，但已经学会了代码的基本格式\protect\footnotemark}
\end{figure}
\footnotetext{Slides 第27页。}

这正是当今大语言模型的雏形——从字符级预测发展到 token 级预测，从 RNN 发展到 Transformer，但\textbf{自回归预测下一个 token 的核心思想始终未变}。

\subsection{本章小结}

字符级语言模型展示了 RNN 在文本生成方面的潜力。其核心思想——自回归地预测下一个 token——直接启发了现代大语言模型的设计。从 Karpathy 2015 年的``The Unreasonable Effectiveness of Recurrent Neural Networks''到 GPT 系列，这一范式延续至今。

%% ========== 隐藏状态可视化 ==========

